% ============================================================
% Sección 3: Metodología
% ============================================================

\section{Metodología}
\label{sec:metodologia}

\subsection{Formulación del problema}
\label{subsec:formulacion}

Se considera el dominio cuadrado adimensional $\Omega = [0,1]^2$, donde las coordenadas
han sido normalizadas por la longitud de onda del láser ($\lambda = 638\,\mathrm{nm}$).
Bajo esta normalización, el número de onda adimensional es $\tilde{k} = 2\pi$.

\paragraph{Validación 1D.}
La ecuación de Helmholtz unidimensional en $x \in [0,1]$ con condición de Dirichlet:
%
\begin{equation}
  \frac{d^2 E}{dx^2} + k^2 E = 0, \quad
  E(x_j) = \cos(k x_j), \quad x_j \in \{0,\ 0.25,\ 0.5,\ 0.75,\ 1\},
  \label{eq:helmholtz1d}
\end{equation}
%
admite la solución analítica $E_{\mathrm{exacta}}(x) = \cos(kx)$, utilizada para
validar la arquitectura SIREN. Se imponen cinco puntos de frontera en lugar de
únicamente los dos extremos: la condición simétrica $E(0) = E(1) = 1$ colapsa
la red hacia la solución constante trivial durante el entrenamiento, ya que
ambos extremos son consistentes con $E \equiv 1$; los tres puntos intermedios
rompen esa simetría y estabilizan la convergencia.

\paragraph{Validación 2D.}
La ecuación de Helmholtz 2D en $\Omega = [0,1]^2$ con condición de onda plana en los
cuatro bordes:
%
\begin{equation}
  E_{\mathrm{exacta}}(x,y) = e^{i(k_x x + k_y y)}, \quad k_x = k_y = \frac{k}{\sqrt{2}},
  \label{eq:solucion2d}
\end{equation}
%
que satisface $k_x^2 + k_y^2 = k^2$.

\subsection{Arquitectura SIREN}
\label{subsec:arquitectura}

Se emplea una arquitectura SIREN con parámetros consolidados a partir de la validación 1D:
$L = 5$ capas ocultas, dimensión $d = 128$ (duplicada respecto a 1D para acomodar el
campo complejo), frecuencia de entrada $\omega_0 = 1.0$ y salida lineal de dos
neuronas $(E_{\mathrm{real}}, E_{\mathrm{imag}})$.
El modelo tiene 66{,}690 parámetros entrenables.

La elección $\omega_0 = 1.0$, en lugar del valor $\omega_0 = 30$ recomendado por
\citet{Sitzmann2020_SIREN} para señales de imagen, está justificada por la naturaleza
del problema: el dominio adimensional $[0,1]^2$ con $k = 2\pi$ no requiere capturar
frecuencias altas, y valores mayores de $\omega_0$ destabilizan el entrenamiento con
L-BFGS.

\subsection{Muestreo por Hipercubo Latino}
\label{subsec:lhs}

En el caso 2D, los $N_c = 3{,}000$ puntos de colocación se generan mediante
Muestreo por Hipercubo Latino (LHS) \citep{Mckay1979_LHS} en lugar de una malla
uniforme.
LHS divide cada dimensión en $N_c$ estratos iguales y coloca exactamente un punto
por estrato, garantizando cobertura uniforme del dominio sin los patrones redundantes
de una malla cartesiana.
Adicionalmente, los $N_b = 300$ puntos de frontera por borde se distribuyen
uniformemente sobre cada lado del dominio.

\subsection{Función de pérdida y pesos}
\label{subsec:loss}

La función de pérdida total para el caso 2D es:
%
\begin{equation}
  \mathcal{L} = \mathcal{L}_{\mathrm{datos}}
              + \lambda_{\mathrm{fís}}\!\left(
                  \mathcal{L}_{R_{\mathrm{real}}} + \mathcal{L}_{R_{\mathrm{imag}}}
                \right),
  \label{eq:loss_2d}
\end{equation}
%
con peso de física $\lambda_{\mathrm{fís}} = 0.1$.
Este valor, calibrado empíricamente, equilibra el ajuste a los datos de frontera y
el cumplimiento de la EDP: valores superiores ($\lambda = 1.0$) sobrepesan la física
e impiden la convergencia en 2D; valores inferiores producen campos que violan la
ecuación de Helmholtz.
La \cref{tab:ablacion_lambda} cuantifica el efecto sobre el error $L^2$ promedio.
Esta fragilidad frente al peso relativo de los términos de pérdida no es un fenómeno
aislado de esta implementación: \citet{Krishnapriyan2021_failureModes} documentan
sistemáticamente que las PINN con pesos fijos entre el término de datos y el término de
residuo físico pueden fallar en converger incluso cuando la pérdida global parece
decrecer, atribuyendo el problema a un paisaje de optimización mal condicionado por la
competencia entre ambos términos.

\begin{table}[htbp]
\centering
\begin{threeparttable}
\caption{Ablación del peso de física $\lambda_{\mathrm{fís}}$ - NB02}
\label{tab:ablacion_lambda}
\begin{tabular}{lccc}
\toprule
$\lambda_{\mathrm{fís}}$ & Error $L^2$ prom.\ (\%) & Converge & Observación \\
\midrule
0.01  & 0.163 & Sí & Física subponderada \\
0.1   & 0.171 & Sí & \textbf{Configuración óptima} \\
1.0   & $>5.0$ & No & Física sobreponderada \\
\bottomrule
\end{tabular}
\begin{tablenotes}
\small
\item Arquitectura SIREN $5\times128$, $\omega_0 = 1.0$, SEED = 42.
      Mismo protocolo Adam + L-BFGS que NB02.
      ``No converge'' indica que el optimizador Adam no alcanzó la condición
      de parada antes de las 15{,}000 épocas con $\varepsilon_{L^2} < 5\%$.
\end{tablenotes}
\end{threeparttable}
\end{table}

\subsection{Estrategia de optimización}
\label{subsec:optimizacion}

El entrenamiento sigue una estrategia bifásica:

\paragraph{Fase 1 - Adam.}
Optimizador Adam con tasa de aprendizaje $\eta = 10^{-3}$, hasta 15{,}000 épocas,
con \textit{early stopping} por paciencia de 800 épocas sin mejora mayor que
$\delta = 10^{-7}$.
Adam realiza exploración global del espacio de parámetros y proporciona un punto
de inicio próximo al mínimo global.

\paragraph{Fase 2 - L-BFGS.}
Optimizador L-BFGS con búsqueda de línea \textit{strong Wolfe}, máximo de 1{,}000
iteraciones y tamaño de historial 100.
L-BFGS aprovecha la información de curvatura del paisaje de pérdida para realizar
un ajuste fino de alta precisión desde el punto entregado por Adam.

Esta combinación, práctica extendida en la literatura de PINNs \citep{Cuomo2022_review},
supera a Adam puro en precisión final y a L-BFGS puro en robustez ante mínimos locales.

\subsection{Métrica de validación}
\label{subsec:metricas}

\paragraph{Error L2 relativo} (NB01 y NB02):
%
\begin{equation}
  \varepsilon_{L2} = \frac{\|\hat{E} - E_{\mathrm{exacta}}\|_2}{\|E_{\mathrm{exacta}}\|_2},
  \label{eq:l2}
\end{equation}
%
evaluado sobre una malla uniforme de $100 \times 100$ puntos.
No existe en la literatura de PINNs un criterio estandarizado de tolerancia de error
para este tipo de validación, por lo que se fijó $\varepsilon_{L2} < 5\%$ como umbral
ingenieril propio de este trabajo. Este valor resulta, además, considerablemente menor
que el orden de magnitud reportado en implementaciones comparables de PINN para la
ecuación de Helmholtz (\cref{tab:comparativa}).

\subsection{Implementación}
\label{subsec:implementacion}

Toda la implementación está en Python con PyTorch, ejecutada en GPU NVIDIA RTX 5050
(8.5 GB VRAM).
Los módulos \texttt{src/models.py}, \texttt{src/losses.py} y \texttt{src/utils.py}
contienen las clases y funciones reutilizables; los notebooks reproducen los experimentos
de forma autocontenida.
La semilla global $\texttt{SEED} = 42$ garantiza reproducibilidad.
